\section{The logarithm of a truncated stationary Gaussian integral}
\label{sec:truncated-stationary-log}
The local small-field comparison uses a normalized integral near the
actual stationary point. Its first correction has order $\beta^{-1}$.
Two details matter: the cutoff is retained when the coupling parameter
vanishes, and the stationary Hessian varies with the coarse field.
We prove the corresponding parameter estimate with both features.

\subsection{The normalized integral}
Let $E$ be the finite-dimensional space of coarse logarithmic
coordinates and let $\mathbb R^d$ be the fine integration space, with
its Euclidean norm. Let $\varphi(x,y)$ be the real action,
$\psi(x,y)$ the additional real exponent, and $\rho(x)$ the stationary
point. Assume
\[
 D_y\varphi(x,\rho(x))=0,\qquad
                   D_y^2\varphi(x,\rho(x))=M(x)>0.
\]
On the parameter domains under consideration assume smoothness of
$\rho$, $\varphi$, $\psi$ and $\tfrac12\log\det M$, with derivatives
through the orders used below. More precisely, a derivative estimate
through order $k$ for the integral uses orders $k+2$ for $\varphi$
and $\rho$, and order $k$ for $\psi$ and the relative log determinant,
on every segment $(x,\rho(x)+tbz)$ with $0\le t\le1$ and $|z|\le D$.
Let their common bounds be $p$. These are bounds on the actual moving
stationary branch and its Hessian, not on a separately chosen quadratic
approximation.

Define the smooth divided difference
\[
 V(b,x,z)=\int_0^1(1-t)
   D_y^2\varphi(x,\rho(x)+tbz)[z,z],dt.
\]
Taylor's formula and stationarity give
\[
 V(b,x,z)=\frac{\varphi(x,\rho(x)+bz)-\varphi(x,\rho(x))}{b^2}
 \quad(b\ne0),\qquad V(0,x,z)=\tfrac12z^TM(x)z.
\]
The integral formula supplies its derivatives also at $b=0$.
Put $s(x)=\tfrac12\log\det M(x)$ and
\begin{align*}
 A(b,x,z)={}&s(x)-s(0)+V(0,0,z)-V(b,x,z)\\
            &+\psi(x,\rho(x))-\psi(x,\rho(x)+bz).
\end{align*}
With $\gamma_0$ the Gaussian of covariance $M(0)^{-1}$, set
\begin{equation}\label{eq:stationary-normalized-integral}
 R(b,x)=\int_{|z|\le D}e^{A(b,x,z)}\,d\gamma_0(z),
                         \qquad L(b,x)=\log R(b,x).
\end{equation}
The normalization is fixed at the base point in the measure; the
varying determinant is explicitly present in $A$.

At $b=0$ the density identity gives exactly
\begin{equation}\label{eq:stationary-zero-mass}
                   R(0,x)=\gamma_x\{|z|\le D\},
                         \qquad \gamma_x=N(0,M(x)^{-1}).
\end{equation}
Indeed the factor $e^{s(x)-s(0)}$ converts the determinant prefactor,
and $V(0,0,z)-V(0,x,z)$ converts the quadratic exponent. The two
$\psi$ terms cancel. The result is the truncated Gaussian mass,
which need not equal one.

There is also an exact parity identity
\begin{equation}\label{eq:stationary-even}
                            R(-b,x)=R(b,x).
\end{equation}
Changing $z$ to $-z$ preserves the ball and the centered Gaussian.
The segment $\rho(x)+t(-b)(-z)$ equals the original segment, and
the two appearances of $z$ in the quadratic evaluation cancel their
signs. Thus $A(-b,x,-z)=A(b,x,z)$, proving the identity before taking
any expansion.

\subsection{Uniform derivatives on the truncated ball}
The chain rule in the integral defining $V$ gives a polynomial bound
$P_k(p,D)$ for all derivatives of $A$ through order $k$, uniformly
for $|z|\le D$. To see the polynomial nature, the map
$(b,x)\mapsto(x,\rho(x)+tbz)$ has derivatives bounded by fixed
polynomials in $1+p+D$; each chain-rule term has at most $k$
factors. Evaluating the two additional slots at $z$ costs at most
$D^2$. Integration against $1-t$ adds no large constant. The
$\psi$ and log-determinant terms have the same type of bound.

Purchase $k\ge1$ (replace a requested order zero by order one here).
Suppose $|b|<v\le1$, $\|x\|<u\le1$, and
$P_k(p,D)\max(u,v)\le1$, enlarging $P_k$ to include its first
order bound. Since $A(0,0,z)=0$, integration along the parameter
segment from zero gives $|A(b,x,z)|\le1$. Derivatives of $e^A$
through order $k$ are consequently bounded by another fixed
polynomial in $P_k$, times $e$. The cutoff ball is fixed in these
parameters, and the reference Gaussian has mass one. This bound is
integrable and uniform on the open parameter domain, so differentiation
under the integral is justified at every order through $k$.

If $\gamma_0\{|z|\le D\}\ge1/2$, then
\begin{equation}\label{eq:stationary-positive-range}
                              e^{-1}/2\le R(b,x)\le e.
\end{equation}
It follows by integrating $e^{-1}\le e^A\le e$ on the ball.
Derivatives of $\log$ are bounded on this fixed positive interval,
so the chain rule gives a common bound for all derivatives of $L$
through order $k$. If $p,D$ have upper exponent zero, this bound
also has upper exponent zero. The smallness condition on $u,v$
holds eventually whenever both have a fixed negative upper exponent.

\subsection{The cutoff error when the dimension grows}
The auxiliary tori have dimension at most $Q(x)$, where $Q$ is bounded
by a fixed polynomial in $1+(\log\beta_F(x))^2$. Suppose
$M(y)^{-1}\le cI$ uniformly on the coarse-field ball, with $c>0$
fixed. Let $P_F(x)$ be the fine monitor and take $D=d_0P_F(x)$,
where $d_0>0$ is fixed. For the actual fine monitor,
$P_F(x)/\log z\to\infty$, and every fixed polynomial in $\log z$
is eventually bounded by $P_F(x)^\epsilon$ for every $\epsilon>0$,
where $z=\beta_F(x)$. These stronger growth properties are used
in the following tail estimate.

Here is a dimension-explicit tail estimate. If a centered Gaussian
has covariance at most $cI$ in dimension at most $P\ge1$, then
for $r,s\ge0$,
\[
 \mathbb P\{|Z|>Pr\}\le2P\exp(-sr+cs^2/2).
\]
If every coordinate has absolute value less than $r$, then
$|Z|^2\le Pr^2\le P^2r^2$. A union bound and the one-coordinate
exponential-moment estimate therefore prove the assertion.
For $P=1+Q(x)$, the monitor growth gives $P\le\sqrt{P_F(x)}$
eventually. Choose $r=d_0\sqrt{P_F(x)}$ and $s=r/c$. Then $Pr\le D$
and
\begin{equation}\label{eq:stationary-tail}
 1-\gamma_y\{|z|\le D\}
        \le2(1+Q(x))\exp\{-d_0^2P_F(x)/(2c)\}.
\end{equation}
The right side has every negative upper exponent. In particular it
is at most $1/2$ eventually, proving the mass condition above.

If $0\le\epsilon\le1/2$ and $1-\epsilon\le r\le1$, integration
of $1/t$ on $[r,1]$ gives $|\log r|\le2\epsilon$. Applying this to
\eqref{eq:stationary-zero-mass} shows that $L(0,y)$ has a value
smaller than every fixed negative power of $\beta_F$, uniformly
on the allowed coarse ball. Derivatives require an additional
argument; a small value by itself is insufficient.

\subsection{Turning a small value into small jets}
We give the finite-order interpolation used for that purpose. Fix
$j\in\mathbb N$, $a>0$, and real $d,m$. There is an exponent $B$
such that, for all sufficiently large $z$, the following implication
holds in any finite-dimensional normed spaces. If on
$\overline B(y,z^{-a})$ a function $F$ has derivatives through order
$j+1$ bounded by $z^d$, and its value is bounded by $z^B$, then its
jets through order $j$ at $y$ are bounded by $z^m$.

The elementary one-step estimate is
\[
             \|D^{r+1}F(y)\|\le 2v/s+Cs
\]
whenever $\|D^rF\|\le v$ and $\|D^{r+2}F\|\le C$ on the
radius-$s$ ball. Subtract the tangent map of $D^rF$ at $y$,
integrate its derivative along a segment of length $s$, and use
the second-derivative bound to control the tangent-map error.
This gives the estimate in every unit direction, hence in operator
norm, without a dimension factor.

Induct on $j$. The case $j=0$ takes $B=m$. In the induction step
choose
\[
 c=a+\max(0,d-m)+1,\qquad b=m-1-c,\qquad s=z^{-c}.
\]
Eventually $2s\le z^{-a}$. Apply the induction hypothesis with target
exponent $b$ on every radius-$s$ ball whose center lies within $s$
of $y$; those balls lie in the original domain. A sufficiently
negative common value exponent $B$ then bounds jets through order
$j-1$ by $z^b$ on that inner ball. The one-step estimate bounds the
order-$j$ derivative at $y$ by
\[
                            2z^{b+c}+z^{d-c}.
\]
Both exponents are strictly less than $m$, and $b<m$ also bounds
all lower derivatives. Increasing the lower threshold on $z$
proves the assertion. This is a finite induction, with all exponents
chosen before the function or dimension is supplied.

Apply it to $F(y)=L(0,y)$ on a microscopic ball of radius
$z^{-1}$ inside the coarse-field domain. The derivative bound of
upper exponent zero is at most $z$ eventually. The tail estimate
provides any value exponent $B$ the induction requests. Hence,
for every fixed real $e$,
\begin{equation}\label{eq:stationary-zero-jets}
                       \|L(0,\cdot)\|_{J^j,y}\le z^e
\end{equation}
eventually, with the stated microscopic margin.

\subsection{The even second-order gain}
Purchase two additional parameter derivatives in the earlier
bounds. For every $r\le j$, the function
$b\mapsto D_y^rL(b,y)$ is even by
\eqref{eq:stationary-even}. Its first derivative at zero vanishes.
If its second derivative is bounded by $C(x)$ for $|b|\le\epsilon$,
Taylor's formula gives
\[
 \|D_y^rL(\epsilon,y)-D_y^rL(0,y)\|
                         \le C(x)\epsilon^2.
\]
The derivative constant $C$ has upper exponent zero. Combining with
\eqref{eq:stationary-zero-jets}, assuming $\beta>0$ and $\epsilon=\beta^{-1/2}<v$, we obtain
\begin{equation}\label{eq:stationary-log-price}
             \|L(\beta^{-1/2},\cdot)\|_{J^j,y}
                              \le C(x)\beta^{-1}+z^e.
\end{equation}
This is the required first-order inverse-coupling price, uniform
for the prescribed auxiliary volumes and coupling band, with the
closed microscopic ball $\overline B(y,z^{-1})$ contained in the
coarse-field parameter domain, provided
the actual stationary branch, segment derivative and Hessian bounds
stated above hold. Identifying this truncated integral with the
full nonlinear push additionally requires its separate exterior
estimate; that estimate is not part of the parity argument.
